\chapter{Basis, Roll and Delivery}\label{ch:m1:basis-roll-and-delivery}

The index is at 6\,000 and its December future trades at 6\,043. A
commentator reads the forty-three points as optimism. They are interest: the
buyer of the future keeps \$300\,000 in the bank for three months instead of
paying for shares, earns about \$3\,185 on it, forgoes about \$1\,050 of
dividends, and the difference, divided by the multiplier, is forty-three
points to the decimal. The premium will melt to nothing by the third Friday of
December whatever the market does. On 20 April 2020 the same mechanism ran
in reverse in crude oil: holders of a future one day from expiry, facing
delivery of barrels they could not store, paid others \$37.63 a barrel to take
the contracts away. A future is tied to its underlying by an arbitrage and by
a delivery procedure, and both have costs. This chapter prices the tie.

\section{Cost of carry and fair value}

\begin{definition}[Basis]\label{def:m1:basis-roll-and-delivery:basis}
The \emph{basis}\index{basis} of a future is the difference between its price
and the spot price of its underlying, here future minus spot. It converges to
zero at expiry, because at expiry the future \emph{is} the underlying.
\end{definition}

\begin{definition}[Cost of carry and fair value]\label{def:m1:basis-roll-and-delivery:carry}
The \emph{cost of carry}\index{cost of carry} of an asset is what it costs to
hold it until a future date: financing, plus storage and insurance for a
commodity, minus the income it pays. The \emph{fair value}\index{fair value}
of a future is the spot price plus the cost of carry to expiry: the futures
price at which buying the asset and selling the future earns exactly the
financing rate.
\end{definition}

\begin{proposition}[Fair value of an index future]\label{prop:m1:basis-roll-and-delivery:fv}
Let $S$ be the index, $r$ the financing rate (simple, actual/360), $T$ the
time to expiry and $d_i$ the dividends, in index points, paid at times $t_i
\le T$. Then
\[
F^* \;=\; S\,(1 + rT) \;-\; \sum_i d_i\,\bigl(1 + r\,(T - t_i)\bigr).
\]
\end{proposition}
\begin{proof}
Borrow $S$, buy the basket, sell one future at $F$. Invest each dividend to
expiry. At expiry deliver the basket's value against the future: the cash is
$F + \sum_i d_i(1 + r(T-t_i)) - S(1+rT)$ whatever the index does. With no
initial outlay and no risk it must be zero.
\end{proof}

\begin{example}[Forty-three points]\label{ex:m1:basis-roll-and-delivery:43}
On 18 September 2026 with $S = 6\,000$, $r = 4.2\%$, expiry on 18 December (91
days) and dividends of 6.0, 9.5 and 5.5 points paid on 15 October, 16
November and 10 December: interest is $6\,000 \times 0.042 \times 91/360 =
63.70$ points, the dividends are worth 21.09 at expiry, and $F^* = 6\,042.61$.
A market price of 6\,047.20 corresponds to an \emph{implied financing
rate}\index{implied financing rate} of 4.50\%: the \omterm{def:m1:delta-one-instruments:swap}{funding spread} of
\cref{ch:m1:delta-one-instruments}, 30 basis points, seen from the other side.
\end{example}

\begin{omfigure}
\begin{tikzpicture}
\begin{axis}[width=11cm,height=4.8cm,
  xlabel={days to expiry},ylabel={fair basis (index points)},x dir=reverse,
  tick label style={font=\small},label style={font=\small},
  xmin=0,xmax=91,ymin=0,ymax=48,grid=major,grid style={gray!25},table/col sep=comma]
\addplot[omDef,very thick,const plot] table[x=days_to_expiry,y=basis]{figdata/markets-1/21-basis-roll-and-delivery/basis.csv};
\node[font=\footnotesize,anchor=west] at (axis cs:63,33) {dividend 6.0};
\node[font=\footnotesize,anchor=west] at (axis cs:31,27) {9.5};
\node[font=\footnotesize,anchor=west] at (axis cs:7.5,12) {5.5};
\end{axis}
\end{tikzpicture}

\omcaption{The fair \omterm{def:m1:basis-roll-and-delivery:basis}{basis} of \cref{ex:m1:basis-roll-and-delivery:43} with the
index held at 6\,000: interest accrues away at 0.7 point a day, and the \omterm{def:m1:basis-roll-and-delivery:basis}{basis}
jumps \emph{up} each time a dividend is paid, since that dividend is no
longer ahead. A \omterm{def:m1:basis-roll-and-delivery:basis}{basis} series has this sawtooth in it before any trading
happens. Data: the chapter's build.}\label{fig:m1:basis-roll-and-delivery:basis}
\end{omfigure}

\section{Index arbitrage}

\begin{proposition}[The arbitrage band]\label{prop:m1:basis-roll-and-delivery:band}
With one-way trading costs $c_S$ on the basket and $c_F$ on the future (as
fractions of value), a financing spread $\phi$ for the arbitrageur when long
the basket and a borrow fee $\beta$ when short it, the future can trade
without arbitrage anywhere in
\[
\Bigl[F^* - S\bigl(2(c_S + c_F) + \beta T\bigr),\;\; F^* + S\bigl(2(c_S + c_F) + \phi T\bigr)\Bigr].
\]
The band is narrower for firms with lower costs, shrinks as expiry
approaches, and is asymmetric when borrowing the basket costs more than
financing it.
\end{proposition}

\begin{example}[A band of seventeen points]\label{ex:m1:basis-roll-and-delivery:band}
With $c_S = 3$ and $c_F = 0.5$ basis points, $\phi = 15$ and $\beta = 40$ basis
points a year, the band of \cref{ex:m1:basis-roll-and-delivery:43} runs from
$F^* - 10.27$ to $F^* + 6.47$: from 6\,032.35 to 6\,049.09. At 6\,050.50 the
future is 7.89 points rich and the trade, buy the basket and sell the future,
locks in 1.41 points, \$70 a contract, held to expiry. Most index arbitrage is
not held to expiry: the position is unwound when the mispricing reverses,
which earns the round trip of the mispricing and pays trading costs twice
more.
\end{example}

\begin{omfigure}
\begin{tikzpicture}
\begin{axis}[width=11cm,height=4.8cm,
  xlabel={time (arbitrary steps)},ylabel={future minus fair value},
  tick label style={font=\small},label style={font=\small},
  xmin=0,xmax=400,ymin=-13,ymax=9,grid=major,grid style={gray!25},table/col sep=comma]
\addplot[omDef,thick] table[x=t,y=mispricing]{figdata/markets-1/21-basis-roll-and-delivery/mispricing.csv};
\addplot[omThm,thick,dashed] table[x=t,y=upper]{figdata/markets-1/21-basis-roll-and-delivery/mispricing.csv};
\addplot[omThm,thick,dashed] table[x=t,y=lower]{figdata/markets-1/21-basis-roll-and-delivery/mispricing.csv};
\end{axis}
\end{tikzpicture}

\omcaption{A simulated mispricing inside the band of
\cref{ex:m1:basis-roll-and-delivery:band} (dashed, in index points). The
upper bound is closer than the lower, because financing a long basket is
cheaper than borrowing a short one, and is reached more often. Data: the
tutorial's simulation.}\label{fig:m1:basis-roll-and-delivery:mispricing}
\end{omfigure}

\section{The roll}

\begin{definition}[Roll]\label{def:m1:basis-roll-and-delivery:roll}
To \emph{roll}\index{roll} a futures position is to close it in the expiring
contract and reopen it in a later one, normally as one calendar-spread trade
(\cref{def:m1:matching-algorithms-and-implied-spreads:calendar}). The price of
the roll is the \omterm{def:m1:matching-algorithms-and-implied-spreads:calendar}{calendar spread}; its \emph{richness} is the spread's excess
over the difference of the two \omterm{def:m1:basis-roll-and-delivery:carry}{fair values}, usually quoted as an annualised
rate.
\end{definition}

For an equity index the \omterm{def:m1:basis-roll-and-delivery:roll}{roll} prices three months of financing less three
months of dividends. A long holder rolling at a spread 4.55 points above \omterm{def:m1:basis-roll-and-delivery:carry}{fair
value} on an index of 6\,000 pays 30 basis points a year more than the
benchmark for its leverage; the short holder on the other side earns them.
The \omterm{def:m1:basis-roll-and-delivery:roll}{roll} takes place over about a week (\cref{fig:m1:futures-contracts-and-exchanges:roll}),
and since everyone must do it, its richness is the cleanest public measure of
the price of balance sheet.

\begin{definition}[Contango and backwardation]\label{def:m1:basis-roll-and-delivery:contango}
A futures curve is in \emph{contango}\index{contango} when later expiries
trade above nearer ones, and in \emph{backwardation}\index{backwardation} when
they trade below.
\end{definition}

For a financial asset the curve's slope is arithmetic: rate minus yield. For
a commodity it includes storage and a \emph{convenience yield}, the value of
having the physical commodity to hand, which is high when inventories are
low. The slope is a return: a long position rolled down a backwardated curve
earns the slope even if the spot price never moves, and one rolled up a
\omterm{def:m1:basis-roll-and-delivery:contango}{contango} curve pays it.

\begin{omfigure}
\begin{tikzpicture}
\begin{axis}[width=11cm,height=4.8cm,
  xlabel={months},ylabel={value of 1 invested},
  tick label style={font=\small},label style={font=\small},
  xmin=0,xmax=24,ymin=0.6,ymax=1.4,grid=major,grid style={gray!25},
  legend columns=2,legend style={font=\footnotesize,at={(0.5,-0.32)},anchor=north,draw=none,fill=none,/tikz/every even column/.append style={column sep=8pt}},
  table/col sep=comma]
\addplot[omThm,very thick] table[x=month,y=contango]{figdata/markets-1/21-basis-roll-and-delivery/rolled.csv};
\addplot[omDef,very thick,dashed] table[x=month,y=backwardation]{figdata/markets-1/21-basis-roll-and-delivery/rolled.csv};
\legend{second month 1.2\% above the first,second month 1.0\% below the first}
\end{axis}
\end{tikzpicture}

\omcaption{A long position rolled monthly for two years while the spot price
never moves from 70. In \omterm{def:m1:basis-roll-and-delivery:contango}{contango} it loses 25\%; in \omterm{def:m1:basis-roll-and-delivery:contango}{backwardation} it gains
27\%. An investor who ``bought oil'' through futures bought the curve as well.
Data: the tutorial's simulation.}\label{fig:m1:basis-roll-and-delivery:rolled}
\end{omfigure}

\section{Delivery and final settlement}

\begin{definition}[Cash settlement and physical delivery]\label{def:m1:basis-roll-and-delivery:delivery}
At expiry a \emph{cash-settled}\index{cash settlement} future pays a last
\omterm{def:m1:clearing-and-settlement:margin}{variation margin} against a final settlement price and disappears. A
\emph{physically delivered}\index{physical delivery} future obliges the short
to deliver, and the long to receive and pay for, the underlying, under the
exchange's delivery rules: grade, location, dates, and, for bond futures, a
choice among deliverable issues.
\end{definition}

\begin{definition}[Special opening quotation]\label{def:m1:basis-roll-and-delivery:soq}
A \emph{special opening quotation}\index{special opening quotation} is a
final settlement value of an index computed from the opening prices of each
of its components on the settlement day. It is not an index level that
anyone saw: the components open at different moments, and the quotation can
lie outside the day's range of the published index.
\end{definition}

\begin{dated}{2026-09}{How three contracts end}\label{dat:m1:basis-roll-and-delivery:ends}
\textbf{E-mini S\&P 500}: trading in the expiring contract stops at the
scheduled opening of the New York Stock Exchange on the third Friday; final
settlement is a special opening quotation from that morning's opening
prices. \textbf{EURO STOXX 50 future}: trading stops at noon on the third
Friday, against the average of the index over the last ten minutes.
\textbf{WTI crude}: \omterm{def:m1:basis-roll-and-delivery:delivery}{physical delivery} at Cushing; trading ends some days
before the delivery month begins, and a holder who is still long then owns
oil in a pipeline hub in Oklahoma.
\end{dated}

\begin{omfigure}
\begin{tikzpicture}[font=\small]
\draw[thick,-{Stealth[length=2.5mm]}] (0,0) -- (13.2,0);
\foreach \x/\d/\t in {1.0/Thursday 16:00/last stock-market close,5.2/Friday 09:30/futures stop trading,11.0/{Friday, later}/quotation published}{
  \draw[thick] (\x,0.12) -- (\x,-0.12);
  \node[above=3pt,font=\footnotesize\bfseries] at (\x,0.12) {\d};
  \node[below=3pt,font=\footnotesize] at (\x,-0.12) {\t};}
\foreach \x in {5.5,5.7,5.8,6.0,6.1,6.4,6.6,7.1,7.9,9.2}{\draw[omThm,thick] (\x,0.22) -- (\x,0);}
\draw[decorate,decoration={brace,amplitude=5pt},omThm,thick] (5.45,0.95) -- (9.25,0.95) node[midway,above=6pt,font=\footnotesize,text=omThm,align=center] {each stock's opening auction prints:\\the quotation is assembled from these};
\draw[decorate,decoration={brace,amplitude=5pt,mirror},omDef,thick] (1.0,-0.95) -- (5.2,-0.95) node[midway,below=6pt,font=\footnotesize,text=omDef,align=center] {overnight: futures trade,\\shares do not};
\end{tikzpicture}

\omcaption{The expiry of a \omterm{def:m1:basis-roll-and-delivery:delivery}{cash-settled} index future that uses a special
opening quotation (New York times). Between the moment futures stop trading
and the moment the last component opens, the holder's exposure is to opening
auctions it can join but no longer hedge with the expiring
contract.}\label{fig:m1:basis-roll-and-delivery:expiry}
\end{omfigure}

An index arbitrageur who is long the basket and short the expiring future
exits by selling every stock in its opening auction on expiry morning: the
proceeds are, by construction, the special opening quotation, and the
position disappears with no \omterm{def:m1:basis-roll-and-delivery:basis}{basis} risk. This is why the third Friday's
opening auctions are among the largest of the quarter
(\cref{ch:m1:auctions}).

\begin{example}[Minus thirty-seven dollars]\label{ex:m1:basis-roll-and-delivery:wti}
The May 2020 WTI future was due to expire on 21 April. On 20 April it opened
the session at \$17.73; between about 14:08 and the end of the settlement
period at 14:30 New York time it traded below zero, for the first time since
the contract was listed in 1983, touched $-\$40.32$ and settled at
$-\$37.63$. Every other expiry settled at a positive price. The regulator's
interim report cites an oversupplied market, the collapse of demand in the
pandemic and concern about storage at the delivery point, and records that
about a fifth of that day's volume in the contract was traded \emph{at
settlement}, at a price to be known only after 14:30. A long who cannot take
delivery must sell before expiry at whatever price clears; the convergence
of future and spot is enforced by those who can deliver and store, and on
that day few could.
\end{example}

\section{Trading at a price not yet known}

\begin{definition}[Exchange for physical]\label{def:m1:basis-roll-and-delivery:efp}
An \emph{exchange for physical}\index{exchange for physical} (EFP) is a
privately negotiated transaction, reported to the exchange, in which one
party buys the cash asset and sells the future and the other does the
opposite, at an agreed price difference. It moves a position between the
cash and futures markets without touching either order book.
\end{definition}

\begin{definition}[Trade at settlement and basis trade at index close]\label{def:m1:basis-roll-and-delivery:tas}
A \emph{trade at settlement}\index{trade at settlement} (TAS) is a futures
trade executed during the session at a price equal to that day's settlement
price, plus or minus an agreed number of ticks, which is known only after the
settlement. A \emph{basis trade at index close}\index{basis trade at index
close} (BTIC) is a futures trade at a price equal to that day's official
closing level of the index plus an agreed \omterm{def:m1:basis-roll-and-delivery:basis}{basis}.
\end{definition}

Both answer the same need: a fund valued at the close wants futures at the
close, or at a fixed distance from it, without trading in the last thirty
seconds. A BTIC quote is a direct quote of the \omterm{def:m1:basis-roll-and-delivery:basis}{basis}, and so of the \omterm{ex:m1:basis-roll-and-delivery:43}{implied
financing rate}: it is where the \omterm{def:m1:delta-one-instruments:swap}{funding spread} of index futures is most
cleanly observed. Both also concentrate risk in the settlement window: the
counterparties of TAS orders hedge in the market before the settlement, and
on 20 April 2020 that hedging pressed on a market with no buyers.

\section{Tutorial: fair value with discrete dividends}\label{tut:m1:basis-roll-and-delivery:fv}

\begin{tutorial}
\textbf{Goal.} Compute the \omterm{def:m1:basis-roll-and-delivery:carry}{fair value} of an index future from dated
dividends, invert it for the \omterm{ex:m1:basis-roll-and-delivery:43}{implied financing rate}, draw the \omterm{def:m1:basis-roll-and-delivery:basis}{basis} over the
contract's life and the arbitrage band. \textbf{End state:} the three data
figures of this chapter.
\begin{tutsteps}
\item \textbf{\omterm{def:m1:basis-roll-and-delivery:carry}{Fair value}}, with each dividend reinvested to expiry.
\omcode{code/firm/fairvalue/firm_fairvalue.py}{21}{27}{Dividends carried to expiry, and the fair value.}\label{lst:m1:basis-roll-and-delivery:fv}
\item \textbf{Implied rate} by bisection, since the dividends' reinvestment
depends on the rate too.
\omcode{code/firm/fairvalue/firm_fairvalue.py}{30}{42}{The financing rate at which the market price is fair.}\label{lst:m1:basis-roll-and-delivery:implied}
\item \textbf{The band}, asymmetric when the basket is dear to borrow.
\omcode{code/firm/fairvalue/firm_fairvalue.py}{53}{62}{Futures prices beyond which index arbitrage pays.}\label{lst:m1:basis-roll-and-delivery:band}
\item \textbf{A rolled position} in a commodity curve.
\omcode{code/markets-1/21-basis-roll-and-delivery/python/basis_demo.py}{46}{53}{Value of a long position rolled each period.}\label{lst:m1:basis-roll-and-delivery:rolled}
\end{tutsteps}
\textbf{What to change next.} Make the dividends uncertain: shift each by
$\pm 10\%$ and see how much of the band that uses up for a one-year future.
This is why long-dated index futures and dividend futures exist as separate
markets (\cref{ch:m1:index-and-single-stock-futures-worldwide}).
\end{tutorial}

\section{Build: the fair-value calculator}\label{bld:m1:basis-roll-and-delivery:fv}

\begin{build}
\bfield{Purpose} The miniature firm's futures \omterm{def:m1:what-a-trading-firm-does:market-maker}{market maker} quotes around
\omterm{def:m1:basis-roll-and-delivery:carry}{fair value}, its index-arbitrage strategy trades against it, and its risk
system uses it to split futures P\&L into index, carry and mispricing.
\bfield{Interface} \texttt{Dividend(pay\_date, points)};
\texttt{fair\_value(spot, rate, today, expiry, divs)};
\texttt{implied\_rate(future, spot, today, expiry, divs)};
\texttt{ArbCosts}; \texttt{arbitrage\_band(\ldots)} returning the lower and
upper no-arbitrage prices; \texttt{roll\_richness\_bp(spread\_market, spot,
rate, today, near, far, divs)}.
\bfield{Rules} Simple interest, actual/360; dividends counted if paid after
today and on or before expiry; expiry dates from the contract master
(\cref{bld:m1:futures-contracts-and-exchanges:master}); dividend points from
the corporate-action adjuster and index weights
(\cref{bld:m1:shares-corporate-actions-indices:adjuster}); borrow fees from
the borrow book (\cref{bld:m1:stock-loan-and-short-selling:borrow}).
\bfield{Acceptance tests} \texttt{code/firm/fairvalue/tests/}: simple
interest with no dividends; the chapter's example by hand; the implied rate
as inverse; the asymmetric band; a \omterm{def:m1:basis-roll-and-delivery:roll}{roll} 30 basis points rich.
\bfield{Stretch} A term structure of rates instead of one $r$; dividends in
the index's currency for a quanto future; an ex-date, not pay-date,
convention, with the difference reported.
\end{build}

\begin{omsources}
\item CME, Rulebook Chapter 358, rules 35802.G and 35803.A (termination of
trading and final settlement), as filed with the CFTC.
\item US Commodity Futures Trading Commission, \emph{Interim Staff Report:
Trading in NYMEX WTI Crude Oil Futures Contract Leading up to, on, and around
April 20, 2020}, November 2020.
\item US Commodity Futures Trading Commission, \emph{Futures Glossary}
(exchange for physicals); CME Group, \emph{Basis Trade at Index Close (BTIC)}.
\item Eurex, \emph{EURO STOXX 50 Index Futures}, product page.
\end{omsources}

\section{Exercises}

\begin{exercise}[$\star$]\label{exo:m1:basis-roll-and-delivery:1}
With $S = 6\,000$, $r = 4.2\%$, 91 days and no dividends, give the \omterm{def:m1:basis-roll-and-delivery:carry}{fair value}
and the \omterm{def:m1:basis-roll-and-delivery:basis}{basis}.
\end{exercise}

\begin{exercise}[$\star$]\label{exo:m1:basis-roll-and-delivery:2}
Add the three dividends of \cref{ex:m1:basis-roll-and-delivery:43}, paid 64,
32 and 8 days before expiry. Give their value at expiry and the \omterm{def:m1:basis-roll-and-delivery:carry}{fair value}.
\end{exercise}

\begin{exercise}[$\star$]\label{exo:m1:basis-roll-and-delivery:3}
A crude-oil curve has its first two monthly contracts at 70.00 and 70.84.
Name the shape and give the annualised \omterm{def:m1:basis-roll-and-delivery:roll}{roll} yield of a long position. Same
for 70.00 and 69.30.
\end{exercise}

\begin{exercise}[$\star\star$]\label{exo:m1:basis-roll-and-delivery:4}
With the costs of \cref{ex:m1:basis-roll-and-delivery:band}, the future trades
at 6\,050.50. Describe the arbitrage, leg by leg, and give its locked-in
profit per contract. What could still go wrong?
\end{exercise}

\begin{exercise}[$\star\star$]\label{exo:m1:basis-roll-and-delivery:5}
The future of \cref{ex:m1:basis-roll-and-delivery:43} trades at 6\,047.20.
Verify the \omterm{ex:m1:basis-roll-and-delivery:43}{implied financing rate}. If the benchmark rises by 25 basis points
tomorrow with the index unchanged, by how much does the \omterm{def:m1:basis-roll-and-delivery:carry}{fair value} change?
\end{exercise}

\begin{exercise}[$\star\star$]\label{exo:m1:basis-roll-and-delivery:6}
On an expiry Friday the published index opens at 5\,990, trades between
5\,985 and 6\,020 during the day, and the special opening quotation is
5\,982. Explain how the quotation can lie below the day's low. Who cares?
\end{exercise}

\begin{exercise}[$\star\star\star$]\label{exo:m1:basis-roll-and-delivery:7}
\emph{Coding.} With \texttt{rolled\_index} reproduce the two end values of
\cref{fig:m1:basis-roll-and-delivery:rolled}. Then let the \omterm{def:m1:basis-roll-and-delivery:contango}{contango} shrink
linearly from 1.2\% to zero over the 24 months and report the end value.
\end{exercise}

\begin{exercise}[$\star\star\star$]\label{exo:m1:basis-roll-and-delivery:8}
\emph{Find the flaw.} ``S\&P futures are trading 43 points above the index
this morning: futures traders expect the market to rise.'' Correct the
statement. What \emph{would} the futures tell you before the stock market
opens?
\end{exercise}

\section{Problem: Rolling a Billion}

\begin{problem}[{Weekend problem --- the quarterly cost of staying long}]\label{pb:m1:basis-roll-and-delivery:1}
A fund holds \$1 billion of S\&P 500 exposure in E-minis (index 6\,000,
multiplier \$50), its cash invested at the benchmark rate of 4.2\%. It is 18
September 2026; it holds the December contract (expiry 18 December) and will
\omterm{def:m1:basis-roll-and-delivery:roll}{roll} to March (19 March 2027). Dividends between the two expiries are 6.0,
9.5 and 5.5 points paid on 15 January, 15 February and 10 March.

\medskip\noindent\textbf{Part I --- The fair \omterm{def:m1:basis-roll-and-delivery:roll}{roll}.}
\begin{enumerate}
\item How many contracts does the fund hold?
\item Give the \omterm{def:m1:basis-roll-and-delivery:carry}{fair value} of the December contract.
\item Give the \omterm{def:m1:basis-roll-and-delivery:carry}{fair value} of the March contract.
\item Give the fair \omterm{def:m1:matching-algorithms-and-implied-spreads:calendar}{calendar spread} (March minus December).
\item Decompose it into interest and dividends.
\end{enumerate}

\medskip\noindent\textbf{Part II --- The market \omterm{def:m1:basis-roll-and-delivery:roll}{roll}.}
The spread trades 4.55 points above fair.
\begin{enumerate}[resume]
\item Give the richness in basis points a year.
\item Give its cost in dollars for this \omterm{def:m1:basis-roll-and-delivery:roll}{roll}.
\item The fund pays half a spread tick (the tick is 0.05) in execution.
Give that cost.
\item Give the total annual cost of rolling, in dollars and in basis points
of the exposure.
\item Compare with the 7 basis points a year of an index fund
(\cref{ch:m1:delta-one-instruments}). What else must enter the comparison?
\end{enumerate}

\medskip\noindent\textbf{Part III --- When to \omterm{def:m1:basis-roll-and-delivery:roll}{roll}.}
\begin{enumerate}[resume]
\item Most \omterm{def:m1:futures-contracts-and-exchanges:oi}{open interest} \omterm{def:m1:basis-roll-and-delivery:roll}{rolls} in the week before expiry. What is the
argument for rolling with the crowd, and for rolling earlier?
\item In December the richness is typically highest. Why?
\item The fund could sell the \omterm{def:m1:basis-roll-and-delivery:basis}{basis} by BTIC instead. What does it learn from
a BTIC quote that it does not learn from the futures price?
\item A bank offers to take the other side of the \omterm{def:m1:basis-roll-and-delivery:roll}{roll} at \omterm{def:m1:basis-roll-and-delivery:carry}{fair value} plus 20
basis points. Why can it, and why would it?
\end{enumerate}

\medskip\noindent\textbf{Part IV --- Judgement.}
\begin{enumerate}[resume]
\item A commodity fund faces the same decision in a curve in steep
\omterm{def:m1:basis-roll-and-delivery:contango}{contango}. What differs?
\item A fund that forgets to \omterm{def:m1:basis-roll-and-delivery:roll}{roll} an E-mini position holds it to expiry.
What happens? And if the contract were WTI?
\item Why is the richness of the \omterm{def:m1:basis-roll-and-delivery:roll}{roll} a measure of something larger than
this fund's costs?
\item Who earns the 30 basis points?
\item State the \emph{named result}: the cost of the \omterm{def:m1:basis-roll-and-delivery:roll}{roll} in basis points a
year.
\item In one sentence: why does the future converge to the index?
\end{enumerate}
\end{problem}

\section{Interview questions}

\begin{interviewq}[$\star$ \iqroles{trader, researcher, bank}]\label{iq:m1:basis-roll-and-delivery:1}
Why does an equity-index future trade at a premium to the index? When would
it trade at a discount?
\end{interviewq}

\begin{interviewq}[$\star$ \iqroles{trader, researcher}]\label{iq:m1:basis-roll-and-delivery:2}
Define \omterm{def:m1:basis-roll-and-delivery:contango}{contango} and \omterm{def:m1:basis-roll-and-delivery:contango}{backwardation}. Does \omterm{def:m1:basis-roll-and-delivery:contango}{contango} mean the market expects
prices to rise?
\end{interviewq}

\begin{interviewq}[$\star\star$ \iqroles{trader, researcher}]\label{iq:m1:basis-roll-and-delivery:3}
Walk me through an index arbitrage trade, including how you get out.
\end{interviewq}

\begin{interviewq}[$\star\star$ \iqroles{researcher, mle}]\label{iq:m1:basis-roll-and-delivery:4}
You are building a continuous futures price series for research. What are
your options at each \omterm{def:m1:basis-roll-and-delivery:roll}{roll}, and what is each good for?
\end{interviewq}

\begin{interviewq}[$\star\star$ \iqroles{trader, bank}]\label{iq:m1:basis-roll-and-delivery:5}
How did a crude-oil future settle at a negative price, and could it happen
to an equity-index future?
\end{interviewq}

\begin{interviewq}[$\star\star\star$ \iqroles{trader, researcher}]\label{iq:m1:basis-roll-and-delivery:6}
The S\&P \omterm{def:m1:basis-roll-and-delivery:roll}{roll} is trading 60 basis points rich into December. How do you
trade it, and what stops you from doing it in unlimited size?
\end{interviewq}
